%
% Documento: Introdução
\setcounter{page}{6}
\chapter{INTRODUÇÃO}\label{chap:introducao}
A Engenharia de Software é uma área da computação voltada ao desenvolvimento e à manutenção de sistemas de software por meio da aplicação de métodos, processos e ferramentas que contribuem para a qualidade do produto desenvolvido. Segundo \citeonline{hirama2011}, a Engenharia de Software fornece os fundamentos e os recursos necessários para organizar as atividades de desenvolvimento, buscando produzir sistemas com qualidade e produtividade. Nesse sentido, sua aplicação é importante para desenvolver soluções que atendam às necessidades dos usuários de maneira organizada e sistemática.

A utilização de sistemas informatizados tem se tornado importante para a organização das atividades profissionais, especialmente em ambientes que lidam diariamente com documentos, registros e informações que precisam ser armazenados e recuperados. No contexto jurídico, os escritórios de advocacia precisam administrar clientes, documentos, processos judiciais e extrajudiciais, prazos e demais atividades relacionadas ao atendimento e ao acompanhamento dos casos. Segundo \citeonline{aleixo2019gestao}, a gestão documental em escritórios de advocacia apresenta desafios relacionados à organização, à guarda e à recuperação dos documentos, evidenciando a importância de práticas adequadas para o gerenciamento dessas informações.

Nesse cenário, a utilização de sistemas de informação pode contribuir para centralizar dados e facilitar o acompanhamento das atividades realizadas. \citeonline{pimenta2024}, ao apresentarem o sistema Brsadv, voltado à gestão de escritórios de advocacia, destacam a contribuição das tecnologias informatizadas para a organização e o gerenciamento de dados nesse setor. Dessa forma, o desenvolvimento de plataformas direcionadas às necessidades dos profissionais do Direito representa uma possibilidade de melhorar a organização das informações e apoiar a execução das atividades administrativas e jurídicas.

Entretanto, a utilização de ferramentas distintas para registrar informações de clientes, documentos, processos e prazos pode dificultar o acompanhamento das atividades quando não existe uma organização integrada desses dados. A falta de centralização pode tornar a consulta às informações menos prática e dificultar a identificação do andamento das demandas e das atividades pendentes. Assim, torna-se relevante desenvolver soluções que reúnam essas informações em um ambiente único e permitam acompanhar as atividades de maneira mais organizada.

Diante desse contexto, este trabalho tem como objetivo desenvolver o JurisFlow, uma plataforma web para a gestão de escritórios de advocacia. A proposta consiste em centralizar informações relacionadas aos usuários, clientes, pastas, casos, processos, diligências, prazos e notificações, possibilitando o acompanhamento das atividades em um único ambiente. A plataforma também contempla a distinção entre demandas pré-processuais e processuais, buscando facilitar a organização e o acompanhamento das informações de acordo com a situação de cada caso.

Para alcançar esse objetivo, serão realizadas atividades de levantamento e definição de requisitos, modelagem do sistema, desenvolvimento das funcionalidades propostas e verificação de seu funcionamento em relação aos requisitos estabelecidos. Dessa maneira, pretende-se contribuir para a organização das informações e das atividades dos escritórios de advocacia por meio de uma solução desenvolvida com base nos princípios da Engenharia de Software.
 
\section{PROBLEMA}
Diante das dificuldades relacionadas à organização e ao acesso às informações na rotina dos escritórios de advocacia, surge a seguinte questão: como o desenvolvimento de uma plataforma web pode contribuir para centralizar e facilitar o gerenciamento de clientes, documentos, casos, processos, diligências e prazos, proporcionando maior organização e praticidade no acompanhamento das atividades do escritório?

\section{OBJETIVOS}
\subsection{OBJETIVO GERAL}
Desenvolver uma plataforma web para a gestão de escritórios de advocacia, denominada JurisFlow, com o propósito de centralizar as informações e auxiliar na organização e no acompanhamento das atividades administrativas e jurídicas.

\subsection{OBJETIVOS ESPECÍFICOS}
    \begin{itemize}
    \item Levantar e definir os requisitos funcionais e não funcionais da plataforma, considerando as necessidades relacionadas à gestão de escritórios de advocacia;
    
    \item Modelar a estrutura e as funcionalidades do sistema por meio de diagramas que representem seus principais elementos e comportamentos;
    
    \item Desenvolver funcionalidades para o gerenciamento de usuários, clientes, pastas, casos, processos e diligências;
    
    \item Implementar recursos para o controle de prazos e o envio de notificações relacionadas às atividades cadastradas;
    
    \item Organizar as informações dos casos conforme sua situação, distinguindo demandas pré-processuais e processuais;
    
    \item Verificar o funcionamento das funcionalidades desenvolvidas por meio de testes, considerando os requisitos estabelecidos para a plataforma.
\end{itemize}

\section{JUSTIFICATIVA}
A rotina de um escritório de advocacia envolve diversas atividades, como o acompanhamento de processos, 
organização de documentos, controle de prazos, cadastro de clientes e gerenciamento das informações dos 
casos. Quando essas atividades não são organizadas de forma adequada, podem ocorrer dificuldades no acesso 
às informações, perda de prazos e aumento do tempo gasto em tarefas administrativas.
Nesse contexto, o uso da tecnologia pode contribuir  para melhorar a organização e facilitar o gerenciamento 
das atividades realizadas pelos escritórios de advocacia. Com isso, este trabalho se justifica pela necessidade de 
desenvolver uma plataforma web, onde a utilização da plataforma pode contribuir para uma melhor organização dos dados, 
facilitar o acesso às informações e reduzir a dependência de processos realizados manualmente.
Dessa forma, o desenvolvimento da plataforma busca aplicar os conhecimentos da Engenharia de Software na 
criação de uma solução voltada para uma necessidade presente na rotina dos escritórios de advocacia, utilizando 
a tecnologia como forma de auxiliar na organização e gestão das atividadesve proporcionando maior praticidade no acesso
e controle das informações.



